\QHeading{某次歌唱比赛的评奖标准如下：}
\par\noindent


\Problem{一等奖：$96\sim100$分；二等奖：$92\sim95$分；三等奖：$87\sim91$分；优秀奖：$80\sim86$分。四年级参加这次歌唱比赛的同学的成绩如下。（单位：分）}
\par

\Problem{\bfseries 四（1）班6名同学的歌唱比赛成绩}
\DataTable{6}{编号 & 评委1 & 评委2 & 评委3 & 评委4 & 平均分\\\hline
01 & 93 & 97 & 92 & 98 & {\color{answerred}95}\\\hline
02 & 86 & 93 & 92 & 97 & {\color{answerred}92}\\\hline
03 & 87 & 92 & 88 & 85 & {\color{answerred}88}\\\hline
04 & 86 & 82 & 88 & 80 & {\color{answerred}84}\\\hline
05 & 90 & 87 & 88 & 95 & {\color{answerred}90}\\\hline
06 & 92 & 93 & 93 & 94 & {\color{answerred}93}\\\hline}
\Problem{\bfseries 四（2）班6名同学的歌唱比赛成绩}
\DataTable{6}{编号 & 评委1 & 评委2 & 评委3 & 评委4 & 平均分\\\hline
01 & 89 & 90 & 92 & 89 & {\color{answerred}90}\\\hline
02 & 93 & 92 & 94 & 93 & {\color{answerred}93}\\\hline
03 & 95 & 97 & 98 & 98 & {\color{answerred}97}\\\hline
04 & 89 & 88 & 90 & 89 & {\color{answerred}89}\\\hline
05 & 85 & 83 & 82 & 82 & {\color{answerred}83}\\\hline
06 & 88 & 90 & 92 & 90 & {\color{answerred}90}\\\hline}
\Problem{① 将上面两个统计表补充完整。② 将比赛成绩整理在下面的统计表中。\bfseries 四年级歌唱比赛获奖情况统计表}
\DataTable{3}{奖项 & 四（1）班 & 四（2）班\\\hline
合计 & {\color{answerred}6} & {\color{answerred}6}\\\hline
一等奖 & {\color{answerred}0} & {\color{answerred}1}\\\hline
二等奖 & {\color{answerred}3} & {\color{answerred}1}\\\hline
三等奖 & {\color{answerred}2} & {\color{answerred}3}\\\hline
优秀奖 & {\color{answerred}1} & {\color{answerred}1}\\\hline}




\QHeading{第1问　补充方法：四（1）班01号}

\Answer{用95分作基数，把多出与不足相互补足。多出量为2、3；不足量为2、3。}

\QHeading{多出量合计}

\Answer{
\Formula{\begin{aligned}[t]&\phantom{={}}2+3\\
&=5\text{（分）}\end{aligned}}}

\QHeading{不足量合计}

\Answer{
\Formula{\begin{aligned}[t]&\phantom{={}}2+3\\
&=5\text{（分）}\end{aligned}}}

\QHeading{计算平均数}

\Answer{
\Formula{\begin{aligned}[t]&\phantom{={}}95+(5-5)\div 4\\
&=95+0\div 4\\
&=95+0\\
&=95\text{（分）}\end{aligned}}\par 答：平均分是95分。}

\QHeading{第1问　补充方法：四（1）班02号}

\Answer{用92分作基数，把多出与不足相互补足。多出量为1、5；不足量为6。}

\QHeading{多出量合计}

\Answer{
\Formula{\begin{aligned}[t]&\phantom{={}}1+5\\
&=6\text{（分）}\end{aligned}}}

\QHeading{计算平均数}

\Answer{
\Formula{\begin{aligned}[t]&\phantom{={}}92+(6-6)\div 4\\
&=92+0\div 4\\
&=92+0\\
&=92\text{（分）}\end{aligned}}\par 答：平均分是92分。}

\QHeading{第1问　补充方法：四（1）班03号}

\Answer{用88分作基数，把多出与不足相互补足。多出量为4；不足量为1、3。}

\QHeading{不足量合计}

\Answer{
\Formula{\begin{aligned}[t]&\phantom{={}}1+3\\
&=4\text{（分）}\end{aligned}}}

\QHeading{计算平均数}

\Answer{
\Formula{\begin{aligned}[t]&\phantom{={}}88+(4-4)\div 4\\
&=88+0\div 4\\
&=88+0\\
&=88\text{（分）}\end{aligned}}\par 答：平均分是88分。}

\QHeading{第1问　补充方法：四（1）班04号}

\Answer{用84分作基数，把多出与不足相互补足。多出量为2、4；不足量为2、4。}

\QHeading{多出量合计}

\Answer{
\Formula{\begin{aligned}[t]&\phantom{={}}2+4\\
&=6\text{（分）}\end{aligned}}}

\QHeading{不足量合计}

\Answer{
\Formula{\begin{aligned}[t]&\phantom{={}}2+4\\
&=6\text{（分）}\end{aligned}}}

\QHeading{计算平均数}

\Answer{
\Formula{\begin{aligned}[t]&\phantom{={}}84+(6-6)\div 4\\
&=84+0\div 4\\
&=84+0\\
&=84\text{（分）}\end{aligned}}\par 答：平均分是84分。}

\QHeading{第1问　补充方法：四（1）班05号}

\Answer{用90分作基数，把多出与不足相互补足。多出量为5；不足量为3、2。}

\QHeading{不足量合计}

\Answer{
\Formula{\begin{aligned}[t]&\phantom{={}}3+2\\
&=5\text{（分）}\end{aligned}}}

\QHeading{计算平均数}

\Answer{
\Formula{\begin{aligned}[t]&\phantom{={}}90+(5-5)\div 4\\
&=90+0\div 4\\
&=90+0\\
&=90\text{（分）}\end{aligned}}\par 答：平均分是90分。}

\QHeading{第1问　补充方法：四（1）班06号}

\Answer{用93分作基数，把多出与不足相互补足。多出量为1；不足量为1。}

\QHeading{计算平均数}

\Answer{
\Formula{\begin{aligned}[t]&\phantom{={}}93+(1-1)\div 4\\
&=93+0\div 4\\
&=93+0\\
&=93\text{（分）}\end{aligned}}\par 答：平均分是93分。}

\QHeading{第1问　补充方法：四（2）班01号}

\Answer{用90分作基数，把多出与不足相互补足。多出量为2；不足量为1、1。}

\QHeading{不足量合计}

\Answer{
\Formula{\begin{aligned}[t]&\phantom{={}}1+1\\
&=2\text{（分）}\end{aligned}}}

\QHeading{计算平均数}

\Answer{
\Formula{\begin{aligned}[t]&\phantom{={}}90+(2-2)\div 4\\
&=90+0\div 4\\
&=90+0\\
&=90\text{（分）}\end{aligned}}\par 答：平均分是90分。}

\QHeading{第1问　补充方法：四（2）班02号}

\Answer{用93分作基数，把多出与不足相互补足。多出量为1；不足量为1。}

\QHeading{计算平均数}

\Answer{
\Formula{\begin{aligned}[t]&\phantom{={}}93+(1-1)\div 4\\
&=93+0\div 4\\
&=93+0\\
&=93\text{（分）}\end{aligned}}\par 答：平均分是93分。}

\QHeading{第1问　补充方法：四（2）班03号}

\Answer{用97分作基数，把多出与不足相互补足。多出量为1、1；不足量为2。}

\QHeading{多出量合计}

\Answer{
\Formula{\begin{aligned}[t]&\phantom{={}}1+1\\
&=2\text{（分）}\end{aligned}}}

\QHeading{计算平均数}

\Answer{
\Formula{\begin{aligned}[t]&\phantom{={}}97+(2-2)\div 4\\
&=97+0\div 4\\
&=97+0\\
&=97\text{（分）}\end{aligned}}\par 答：平均分是97分。}

\QHeading{第1问　补充方法：四（2）班04号}

\Answer{用89分作基数，把多出与不足相互补足。多出量为1；不足量为1。}

\QHeading{计算平均数}

\Answer{
\Formula{\begin{aligned}[t]&\phantom{={}}89+(1-1)\div 4\\
&=89+0\div 4\\
&=89+0\\
&=89\text{（分）}\end{aligned}}\par 答：平均分是89分。}

\QHeading{第1问　补充方法：四（2）班05号}

\Answer{用83分作基数，把多出与不足相互补足。多出量为2；不足量为1、1。}

\QHeading{不足量合计}

\Answer{
\Formula{\begin{aligned}[t]&\phantom{={}}1+1\\
&=2\text{（分）}\end{aligned}}}

\QHeading{计算平均数}

\Answer{
\Formula{\begin{aligned}[t]&\phantom{={}}83+(2-2)\div 4\\
&=83+0\div 4\\
&=83+0\\
&=83\text{（分）}\end{aligned}}\par 答：平均分是83分。}

\QHeading{第1问　补充方法：四（2）班06号}

\Answer{用90分作基数，把多出与不足相互补足。多出量为2；不足量为2。}

\QHeading{计算平均数}

\Answer{
\Formula{\begin{aligned}[t]&\phantom{={}}90+(2-2)\div 4\\
&=90+0\div 4\\
&=90+0\\
&=90\text{（分）}\end{aligned}}\par 答：平均分是90分。}

\QHeading{第2问　另一种分类方法}

\Answer{先把12位同学各自的平均分求好，再将平均分从高到低排列，按原题的奖项区间逐个归类。不要把单个评委的评分直接归入奖项，也不要把4个评委的评分当成4名学生。表格结果与前面的填表答案相同。}

\Answer{
\QHeading{计算过程：1. 补全平均分}把每位同学的4个评委分数相加，再除以4。原题没有去掉最高、最低分的要求，因此4个分数均计入。

\subsubsection*{四（1）班}

\par\medskip\noindent\textbf{01号：}\par\nopagebreak\Formula{\begin{aligned}[t]&\phantom{={}}(93+97+92+98)\div4\\&=(190+92+98)\div4\\&=(282+98)\div4\\&=380\div4\\&=95\text{（分）}\end{aligned}}\par\medskip

\par\medskip\noindent\textbf{02号：}\par\nopagebreak\Formula{\begin{aligned}[t]&\phantom{={}}(86+93+92+97)\div4\\&=(179+92+97)\div4\\&=(271+97)\div4\\&=368\div4\\&=92\text{（分）}\end{aligned}}\par\medskip

\par\medskip\noindent\textbf{03号：}\par\nopagebreak\Formula{\begin{aligned}[t]&\phantom{={}}(87+92+88+85)\div4\\&=(179+88+85)\div4\\&=(267+85)\div4\\&=352\div4\\&=88\text{（分）}\end{aligned}}\par\medskip

\par\medskip\noindent\textbf{04号：}\par\nopagebreak\Formula{\begin{aligned}[t]&\phantom{={}}(86+82+88+80)\div4\\&=(168+88+80)\div4\\&=(256+80)\div4\\&=336\div4\\&=84\text{（分）}\end{aligned}}\par\medskip

\par\medskip\noindent\textbf{05号：}\par\nopagebreak\Formula{\begin{aligned}[t]&\phantom{={}}(90+87+88+95)\div4\\&=(177+88+95)\div4\\&=(265+95)\div4\\&=360\div4\\&=90\text{（分）}\end{aligned}}\par\medskip

\par\medskip\noindent\textbf{06号：}\par\nopagebreak\Formula{\begin{aligned}[t]&\phantom{={}}(92+93+93+94)\div4\\&=(185+93+94)\div4\\&=(278+94)\div4\\&=372\div4\\&=93\text{（分）}\end{aligned}}\par\medskip

\subsubsection*{四（2）班}

\par\medskip\noindent\textbf{01号：}\par\nopagebreak\Formula{\begin{aligned}[t]&\phantom{={}}(89+90+92+89)\div4\\&=(179+92+89)\div4\\&=(271+89)\div4\\&=360\div4\\&=90\text{（分）}\end{aligned}}\par\medskip

\par\medskip\noindent\textbf{02号：}\par\nopagebreak\Formula{\begin{aligned}[t]&\phantom{={}}(93+92+94+93)\div4\\&=(185+94+93)\div4\\&=(279+93)\div4\\&=372\div4\\&=93\text{（分）}\end{aligned}}\par\medskip

\par\medskip\noindent\textbf{03号：}\par\nopagebreak\Formula{\begin{aligned}[t]&\phantom{={}}(95+97+98+98)\div4\\&=(192+98+98)\div4\\&=(290+98)\div4\\&=388\div4\\&=97\text{（分）}\end{aligned}}\par\medskip

\par\medskip\noindent\textbf{04号：}\par\nopagebreak\Formula{\begin{aligned}[t]&\phantom{={}}(89+88+90+89)\div4\\&=(177+90+89)\div4\\&=(267+89)\div4\\&=356\div4\\&=89\text{（分）}\end{aligned}}\par\medskip

\par\medskip\noindent\textbf{05号：}\par\nopagebreak\Formula{\begin{aligned}[t]&\phantom{={}}(85+83+82+82)\div4\\&=(168+82+82)\div4\\&=(250+82)\div4\\&=332\div4\\&=83\text{（分）}\end{aligned}}\par\medskip

\par\medskip\noindent\textbf{06号：}\par\nopagebreak\Formula{\begin{aligned}[t]&\phantom{={}}(88+90+92+90)\div4\\&=(178+92+90)\div4\\&=(270+90)\div4\\&=360\div4\\&=90\text{（分）}\end{aligned}}\par\medskip

\par\medskip\noindent
\QHeading{计算过程：2. 按奖项分类统计}一等奖96--100分；二等奖92--95分；三等奖87--91分；优秀奖80--86分。

四（1）班：01、02、06号获二等奖，共3人；03、05号获三等奖，共2人；04号获优秀奖，共1人；一等奖0人。合计
\Formula{\begin{aligned}[t]&\phantom{={}}0+3+2+1\\&=3+2+1\\&=5+1\\&=6\text{（人）}\end{aligned}}
\par 四（2）班：03号获一等奖；02号获二等奖；01、04、06号获三等奖；05号获优秀奖。人数依次为1、1、3、1，合计
\Formula{\begin{aligned}[t]&\phantom{={}}1+1+3+1\\&=2+3+1\\&=5+1\\&=6\text{（人）}\end{aligned}}
\par

}
